\documentclass{article}

\usepackage[backend=biber,sorting=none]{biblatex}
\addbibresource{mybib.bib}

\usepackage{amsmath, amssymb}
\usepackage{graphicx}

\newcommand{\N}{\mathbb{N}}
\newcommand{\iidsim}{\overset{\text{iid}}{\sim}}
\newcommand{\perpci}{\perp\!\!\!\perp}

\title{Per-token-optimal Coupling Dominates Seed Reuse for Coupling Response Distributions}
\author{Malcolm Sharpe}
\date{\today}

\begin{document}

\maketitle

\begin{abstract}
When qualitatively comparing two LLM response distributions, small differences are confounded by the stochasticity of generation. To control for this confound, a common approach is to reuse a seed. An alternative is to optimally couple the next-token distributions, inspired by speculative decoding. We prove that the optimal coupling method, compared to a broad class of seed reuse methods, achieves an equal or greater probability of agreement for every prefix value.
\end{abstract}

\section*{Introduction}

Some LLM APIs provide a seed parameter \cite{sharpe2024blogseedparam}.
Although perhaps originally intended for reproducibility across calls with equal inputs,
it's also useful for controlling stochasticity for calls with different inputs.
When used to couple two response distributions,
we will call this per-token common random numbers (PTCRN).

Standard speculative decoding \cite{leviathan2023fast} \cite{chen2023acceleratinglargelanguagemodel} implicitly defines a coupling of the response distributions of the draft model and target model.
Though in that context the goal is to reduce generation latency,
the same coupling can be applied for controlling stochasticity.
Call this per-token optimal coupling (PTOC).

The contribution of this paper is to prove that PTOC dominates every PTCRN scheme in a particular sense:
every prefix value has at least as high a probability of agreement with PTOC as it does with PTCRN.
As a corollary, PTCRN is also more effective at maximizing the length of the longest common prefix.

\section*{Definitions}

\subsection*{Seed Reuse}

The behavior of seed reuse depends on the implementation of the response generation algorithm. As an example, huggingface's transformers \cite{wolf2020huggingfacestransformersstateoftheartnatural} uses \texttt{torch.\allowbreak multinomial} to sample from each next-token distribution in turn. When the seed is reused, it causes the corresponding calls to \texttt{torch.multinomial} to use the same random numbers.

Inspired by that example implementation, the dependence of response generation on seed can be broadly modeled as follows:

Let $\Sigma$ be a finite set of tokens. Let the response $X \in \Sigma^\N$ be given by, for all $i \geq 0$,

\begin{itemize}
    \item $X_i = f_X(U_i, X_{<i})$
\end{itemize}

where

\begin{itemize}
    \item $U_i \iidsim U[0,1)$
    \item $f_X : [0,1) \times \Sigma^* \to \Sigma$
\end{itemize}

Then seed reuse simply uses the same random variables $U_i$ for two different token choice functions $f_X$ and $f_Y$. This puts it in the family of ``common random numbers'' schemes, as they are known in the variance reduction literature \cite{rubinstein2016simulation}, which motivates describing seed reuse as per-token common random numbers (PTCRN).

\subsection*{Per-token Optimal Coupling (PTOC)}

Standard speculative decoding \cite{leviathan2023fast} \cite{chen2023acceleratinglargelanguagemodel} optimally couples each pair of next-token distributions.

The probability of agreement on a given next token can thus be calculated as follows, a well-known fact from optimal coupling \cite{angel2021pairwiseoptimalcouplingmultiple}. When using PTOC, for every $i \geq 0$ and $a \in \Sigma^\N$,

\begin{equation}
\begin{split}
&P[X_i = Y_i = a_i \mid X_{<i} = Y_{<i} = a_{<i}] \\
&\quad= \min \{ P[X_i = a_i \mid X_{<i} = a_{<i}], P[Y_i = a_i \mid Y_{<i} = a_{<i}] \}
\end{split}
\end{equation}

\section*{Domination Theorem}

\subsection*{Statement}

\textbf{Theorem.} Assume $X$ and $Y$ are coupled using a PTCRN scheme. Then, for every $i \geq 0$ and $a \in \Sigma^\N$,

\begin{equation}
\begin{split}
&P[X_{<i} = Y_{<i} = a_{<i}] \\
&\quad\leq \prod_{j < i} \min \{ P[X_j = a_j \mid X_{<j} = a_{<j}], P[Y_j = a_j \mid Y_{<j} = a_{<j}] \}
\end{split}
\end{equation}

Observe that, for PTOC, this inequality holds with equality.

\subsection*{Conditional Independence Property}

Every PTCRN scheme satisfies these two symmetric conditional independence properties: for all $i \geq 0$,

\begin{itemize}
    \item $X_i \perpci Y_{<i} \mid X_{<i}$
    \item $Y_i \perpci X_{<i} \mid Y_{<i}$
\end{itemize}

The reason is that the value of $Y_{<i}$ only gives information about $U_{<i}$, not $U_i$, and $X_{<i}$ already contains all the relevant information about $U_{<i}$.

\subsection*{Proof}

We are now ready to show that PTOC dominates PTCRN.

\textbf{Proof.}
\begin{align*}
&P[X_{<i} = Y_{<i} = a_{<i}] \\
&\quad= \prod_{j < i} P[X_j = Y_j = a_j \mid X_{<j} = Y_{<j} = a_{<j}] \\
&\quad\leq \prod_{j < i} \min \{ P[X_j = a_j \mid X_{<j} = Y_{<j} = a_{<j}], P[Y_j = a_j \mid X_{<j} = Y_{<j} = a_{<j}] \} \\
&\quad= \prod_{j < i} \min \{ P[X_j = a_j \mid X_{<j} = a_{<j}], P[Y_j = a_j \mid Y_{<j} = a_{<j}] \}
\end{align*}

The reasons for each line:
\begin{itemize}
    \item Auto-regressive factorization.
    \item Optimal coupling of two discrete distributions.
    \item Conditional independence.
\end{itemize}
\hfill \(\square\)

\section*{Related Work}

It is possible to outperform PTOC at maximizing the probability of agreement of every prefix value,
by breaking the conditional independence property with a more sophisticated coupling method \cite{sun2024blockverificationacceleratesspeculative_v1} \cite{sharpe2024blogapoc}.

\section*{Discussion}

The superiority of PTOC over PTCRN only applies when coupling two response distributions. For three or more response distributions, PTOC's performance is no longer achievable, an instance of a general fact for optimal coupling \cite{angel2021pairwiseoptimalcouplingmultiple}.

\printbibliography

\end{document}